%==============================================================================
%  A Thermodynamically-Consistent MeshGraphNet Surrogate for Transient
%  Welding Heat Transfer
%
%  Self-contained LaTeX source for Overleaf. Standard packages only.
%  Figures are guarded with \IfFileExists so the document compiles whether or
%  not you upload the PNGs. To show the figures, upload these files (from
%  data/output/) next to this .tex (or into a ./figs subfolder and adjust the
%  \graphicspath below):
%     spiral_pf_best_ep60_thermal_history.png
%     ablation_figure.png
%     compare_spiral_fem_on_off.png
%     compare_val_sim001_fem_on_off.png
%     compare_val_sim006_fem_on_off.png
%==============================================================================
\documentclass[11pt]{article}

\usepackage[utf8]{inputenc}
\usepackage[T1]{fontenc}
\usepackage{lmodern}
\usepackage[margin=1in]{geometry}
\usepackage{amsmath,amssymb,amsthm,mathtools}
\usepackage{bm}
\usepackage{graphicx}
\usepackage{booktabs}
\usepackage{array}
\usepackage{multirow}
\usepackage{siunitx}
\usepackage{xcolor}
\usepackage{enumitem}
\usepackage{algorithm}
\usepackage{algpseudocode}
\usepackage[hidelinks]{hyperref}
\usepackage[capitalize]{cleveref}  % must load AFTER hyperref

\graphicspath{{./}{./figs/}}

% --- Convenience macros -------------------------------------------------------
\newcommand{\dd}{\mathrm{d}}
\newcommand{\T}{\mathsf{T}}                 % transpose
\newcommand{\grad}{\nabla}
\newcommand{\Lw}{\bm{L}_{\!w}}
\newcommand{\bmu}{\bm{\mu}}
\newcommand{\cpapp}{c_p^{\mathrm{app}}}
\newcommand{\cpsens}{c_p^{\mathrm{sens}}}
\newcommand{\Tinf}{T_{\infty}}
\newcommand{\Rset}{\mathbb{R}}
\newcommand{\softplus}{\operatorname{softplus}}
\newcommand{\norm}[1]{\left\lVert#1\right\rVert}
\newcommand{\abs}[1]{\left\lvert#1\right\rvert}

\theoremstyle{definition}
\newtheorem{proposition}{Proposition}

% --- Figure helper: only include if the file is present -----------------------
\newcommand{\figIfExists}[4]{%
  % #1 filename, #2 width, #3 caption, #4 label
  \begin{figure}[t]
    \centering
    \IfFileExists{#1}{\includegraphics[width=#2\linewidth]{#1}}{%
      \fbox{\parbox[c][0.30\linewidth][c]{#2\linewidth}{\centering\ttfamily
        [figure placeholder]\\ upload \texttt{\detokenize{#1}}}}}
    \caption{#3}
    \label{#4}
  \end{figure}}

\title{\bfseries A Thermodynamically-Consistent MeshGraphNet Surrogate\\[2pt]
       for Transient Welding Heat Transfer:\\[2pt]
       Enthalpy-State GENERIC Structure and Push-Forward Training}

\author{}
\date{}

\begin{document}
\maketitle

\begin{abstract}
\noindent
We develop a graph-neural-network surrogate for transient thermal welding
simulation that replaces an expensive nonlinear finite-element (FEM) solver with
a fast, autoregressive learned model. The surrogate follows the
Encoder--Processor--Decoder \emph{MeshGraphNet} (MPNN) architecture and is trained
on snapshot pairs produced by a 2D transient heat-conduction solver with a moving
Goldak double-ellipsoid arc source, temperature-dependent apparent heat capacity
(latent heats of fusion and vaporization), and mixed convective/radiative/Dirichlet
boundaries. We identify and correct a subtle but consequential
thermodynamic inconsistency that arises when one imposes GENERIC
(\emph{General Equation for the Non-Equilibrium Reversible--Irreversible Coupling})
structure directly on temperature: under strong latent heat, conserving the sum of
nodal temperatures is \emph{not} conserving energy. The fix is the classical
\emph{enthalpy method}: carry the volumetric enthalpy as the conserved state, evolve
it with a learned symmetric positive semi-definite (SPSD) graph-Laplacian dissipation
operator, and recover temperature by inverting the latent-heat-aware enthalpy curve.
Taking the \emph{extensive} nodal energy as the conserved state makes this guarantee
exact on arbitrary (including non-uniform) meshes, and we treat the arc source and
boundary losses through the open-system extension of GENERIC. This guarantees the first
and second laws \emph{by construction}, even through phase change. A learned, per-node \emph{enriched source} restores the melt-pool peak heating
that the structure-preserving dissipation alone cannot supply, without breaking the
guarantees. Finally, we align the training objective with long autoregressive
rollouts using \emph{push-forward} ($K{=}2$) multi-step backpropagation. On a
held-out, out-of-distribution Archimedean-spiral weld (\num{1526} autoregressive
steps), the model attains a global RMSE of \SI{33.2}{K} while respecting the maximum
principle (no sub-ambient cooling) and remaining stable. A matched ablation
(GENERIC enthalpy structure ON vs.\ OFF, both at $K{=}2$) shows that the thermodynamic
structure provides $\sim$$3\times$ lower rollout-RMSE variance, maximum-principle
respect, and a validation metric that tracks the long rollout, whereas the
unconstrained network is accurate-but-fragile.
\end{abstract}

\tableofcontents

%==============================================================================
\section{Introduction}
%==============================================================================

Arc welding involves a moving, highly localized heat source that drives a metal
plate through melting and partial vaporization, leaving steep transient thermal
gradients in its wake. Accurate prediction of the temperature history is essential
for estimating residual stresses, distortion, and microstructure. The standard tool
is the transient nonlinear finite-element (FEM) method, but a single high-fidelity
simulation can take minutes to hours, making design-space exploration,
uncertainty quantification, and online control prohibitively expensive.

\emph{Scientific machine learning} (SciML) offers a route to fast surrogates that
learn the simulator's input--output map. Among mesh-based learned simulators,
\emph{MeshGraphNets}~\cite{pfaff2021learning} are particularly attractive: they
operate directly on the simulation mesh as a graph, generalize across mesh
geometries, and roll forward autoregressively. However, two challenges are acute
for welding:

\begin{enumerate}[leftmargin=1.4em]
  \item \textbf{Long-horizon stability.} Autoregressive surrogates accumulate
  error step over step. A model that is accurate on short validation rollouts can
  diverge catastrophically (spurious far-field heating, sub-ambient cooling, or
  global runaway) over the much longer inference horizons of interest.
  \item \textbf{Phase-change physics.} Welding energy storage is dominated by
  \emph{latent heat}: melting and boiling absorb large amounts of energy at nearly
  constant temperature, so the energy--temperature relationship is strongly
  nonlinear. A surrogate that implicitly assumes energy $\propto$ temperature
  mis-handles exactly the melt pool, which is the region of greatest interest.
\end{enumerate}

This paper addresses both. We build a configurable MeshGraphNet surrogate and equip
it with an optional \emph{structure-preserving thermodynamic head} derived from the
GENERIC framework~\cite{grmela1997dynamics,ottinger2005beyond}. Our contributions are:

\begin{enumerate}[leftmargin=1.4em]
  \item A complete, reproducible pipeline: a nonlinear 2D transient welding FEM
  solver (Goldak source, apparent heat capacity, Robin/Dirichlet boundaries,
  $\theta$-method with Picard iteration), a physics-aware graph representation with
  purely \emph{relative} features, and an Encoder--Processor--Decoder MPNN
  (\cref{sec:fem,sec:graph,sec:mpnn}).
  \item A diagnosis of the thermodynamic inconsistency of imposing GENERIC on
  temperature under latent heat, and its correction via an \emph{enthalpy-state}
  GENERIC head that conserves energy and produces entropy by construction, with a
  numerically well-conditioned temperature-equivalent rescaling
  (\cref{sec:generic,sec:enthalpy}). Choosing the \emph{extensive} nodal energy
  $\mathcal E_i=h_iV_i$ as the state makes energy conservation exact on arbitrary
  (non-uniform) meshes, and the arc/boundary terms are handled by the open-system
  extension of GENERIC.
  \item An \emph{enriched source} that restores melt-pool peak heating while leaving
  the dissipative operator structure-preserving (\cref{sec:enriched}).
  \item A \emph{push-forward} multi-step training scheme that aligns the objective
  with long rollouts (\cref{sec:training}), and a matched ablation isolating the value
  of the thermodynamic structure (\cref{sec:results}).
\end{enumerate}

%==============================================================================
\section{Governing physics and the finite-element solver}
\label{sec:fem}
%==============================================================================

\subsection{Continuum model}

We model the top-down (in-plane) surface of a thin plate of thickness $H_{\!z}$,
and solve the through-thickness-integrated transient heat-conduction equation for
the temperature field $T(\bm{x},t)$, $\bm{x}\in\Omega\subset\Rset^2$:
\begin{equation}
  \rho\,\cpapp(T)\,\frac{\partial T}{\partial t}
  \;=\;
  \grad\!\cdot\!\big(k(T)\,\grad T\big) \;+\; Q(\bm{x},t),
  \label{eq:heat}
\end{equation}
where $\rho$ is density, $k$ thermal conductivity, $Q$ the volumetric heat source,
and $\cpapp(T)$ the \emph{apparent} heat capacity, which folds latent heat into an
effective specific heat (\cref{sec:latent}). Boundaries $\partial\Omega$ carry either
a Dirichlet condition $T=T_D$ or a Robin condition combining convection and
(optionally linearized) radiation,
\begin{equation}
  -k\,\grad T \cdot \bm{n}
  \;=\;
  h_{\mathrm{conv}}\,(T-\Tinf)
  \;+\; \varepsilon\,\sigma_{\mathrm{SB}}\,(T^4-\Tinf^4),
  \label{eq:robin}
\end{equation}
with outward normal $\bm{n}$, convection coefficient $h_{\mathrm{conv}}$, ambient
$\Tinf$, emissivity $\varepsilon$, and Stefan--Boltzmann constant
$\sigma_{\mathrm{SB}}=\SI{5.670e-8}{W.m^{-2}.K^{-4}}$.

\subsection{Latent heat and the enthalpy curve}
\label{sec:latent}

The defining physical feature of welding is that the metal melts (and locally boils):
a phase change absorbs a large \emph{latent heat} at nearly constant temperature, so
$\cpapp$ is strongly temperature-dependent. We model it as a sensible heat capacity
plus two normalized Gaussian bumps centered on melting and boiling:
\begin{equation}
  \cpapp(T) \;=\; \cpsens
  \;+\; L_f\,\frac{\dd f_\ell}{\dd T}
  \;+\; L_v\,\frac{\dd f_v}{\dd T},
  \label{eq:cpapp}
\end{equation}
\begin{equation}
  \frac{\dd f_\ell}{\dd T}=\frac{1}{s_\ell\sqrt{\pi}}\,
    e^{-\left((T-T_m)/s_\ell\right)^2},
  \qquad
  \frac{\dd f_v}{\dd T}=\frac{1}{s_v\sqrt{\pi}}\,
    e^{-\left((T-T_b)/s_v\right)^2},
  \label{eq:bumps}
\end{equation}
where $L_f, L_v$ are the latent heats of fusion and vaporization,
$T_m=\tfrac12(T_{\mathrm{sol}}+T_{\mathrm{liq}})$ is the mean melting temperature,
$T_b$ the boiling point, and $s_\ell, s_v$ smoothing widths. Because each bump
integrates to one, $\int \cpapp\,\dd T$ recovers exactly $L_f$ and $L_v$. The
vaporization bump acts as a physical \emph{thermostat}: energy that would otherwise
superheat near-pool nodes to unphysical values ($\sim$\SI{e4}{K}) is instead
absorbed by the (smeared) phase change, capping the dynamic range near boiling.

Integrating \cref{eq:cpapp} gives the \emph{volumetric enthalpy} (energy stored per
unit volume relative to a reference $T_0$):
\begin{equation}
  h(T) \;=\; \rho \int_{T_0}^{T} \cpapp(\tau)\,\dd\tau .
  \label{eq:enthalpy}
\end{equation}
Away from phase change $h(T)\approx \rho\,\cpsens\,(T-T_0)$ is nearly linear; across
melting/boiling the curve develops plateaus where a large change in $h$ produces an
almost-zero change in $T$ (equivalently the slope $\dd h/\dd T=\rho\,\cpapp$ spikes).
This curve, illustrated schematically in \cref{fig:hcurve}, is the crux of the
thermodynamic argument in \cref{sec:generic}. Numerical parameters
(mild steel) are listed in \cref{tab:material}.

\begin{figure}[t]
\centering
\fbox{\parbox{0.86\linewidth}{\centering\small
\textit{Enthalpy--temperature curve $h(T)$ (schematic).}\\[4pt]
\ttfamily
\begin{tabular}{l}
 h(T)\hfill .--- boiling plateau ($L_v$ absorbed)\\
 \ \ \textbar\hfill \_\_\_/\\
 \ \ \textbar\hfill \_\_\_\_\_/\\
 \ \ \textbar\ \ melting plateau \_\_\_---/\ \ ($L_f$ at $\sim T_m$)\\
 \ \ \textbar\hfill \_\_\_---\\
 \ \ \textbar\ \ \_\_\_-/\ \ (slope $\rho c_p$ away from phase change)\\
 \ \ +\rule{4.2cm}{0.4pt}\textgreater\ T\\
\end{tabular}}}
\caption{The enthalpy--temperature relation $h(T)$. Near melting and boiling a
large change in enthalpy (energy) yields almost no change in temperature: the
phase change absorbs the energy. This nonlinearity is what a temperature-state
GENERIC head mis-models and what the enthalpy-state head captures exactly.}
\label{fig:hcurve}
\end{figure}

\begin{table}[t]
\centering
\caption{Material and source parameters (mild-steel defaults, SI/kelvin).}
\label{tab:material}
\small
\begin{tabular}{lll}
\toprule
Symbol & Quantity & Value \\
\midrule
$\rho$              & density                       & \SI{7850}{kg.m^{-3}} \\
$\cpsens$           & sensible specific heat        & \SI{600}{J.kg^{-1}.K^{-1}} \\
$k$                 & thermal conductivity          & \SI{30}{W.m^{-1}.K^{-1}} \\
$L_f$               & latent heat of fusion         & \SI{2.7e5}{J.kg^{-1}} \\
$T_{\mathrm{sol}},T_{\mathrm{liq}}$ & solidus, liquidus & \SI{1723}{K}, \SI{1773}{K} \\
$L_v$               & latent heat of vaporization   & \SI{6.3e6}{J.kg^{-1}} \\
$T_b$               & boiling point                 & \SI{3134}{K} \\
$s_v$ (\texttt{vap\_width}) & vaporization smoothing & \SI{120}{K} \\
$H_{\!z}$           & plate thickness               & \SIrange{4}{8}{mm} \\
$P$                 & gross arc power               & \SIrange{1500}{3500}{W} \\
$\eta$              & arc efficiency                & \numrange{0.70}{0.90} \\
\bottomrule
\end{tabular}
\end{table}

\subsection{Goldak double-ellipsoid heat source}

The arc is modeled by the Goldak double-ellipsoid distribution~\cite{goldak1984new}.
For the thickness-integrated 2D model we analytically integrate the 3D ellipsoid
through the plate thickness to obtain a surface flux, then divide by $H_{\!z}$ to get
a volumetric source. With the torch at position $\bm{p}(t)$ moving along trajectory
$\bm\gamma(s)$ with unit tangent $\hat{\bm t}$ and normal $\hat{\bm n}$, define the
co-moving coordinates $\xi=(\bm x-\bm p)\cdot\hat{\bm t}$ (along travel) and
$\eta=(\bm x-\bm p)\cdot\hat{\bm n}$ (transverse). Then
\begin{equation}
  Q(\bm x,t) \;=\; \frac{1}{H_{\!z}}\,
  \frac{6\,f\,\eta_{\mathrm{eff}}P}{\pi\,b\,c}\,
  \exp\!\Big(-\frac{3\xi^2}{c^2}-\frac{3\eta^2}{b^2}\Big),
  \label{eq:goldak}
\end{equation}
where $\eta_{\mathrm{eff}}P$ is the net deposited power ($\eta_{\mathrm{eff}}$ the
arc efficiency), $b$ is the transverse half-width, and the along-travel semi-axis and
fraction switch between the front lobe ($c=c_f$, $f=f_f$ for $\xi\ge0$) and the rear
lobe ($c=c_r$, $f=f_r$ for $\xi<0$), with the Goldak constraint $f_f+f_r=2$. During a
post-weld cooling phase the source is switched off ($Q\equiv0$) and the torch is
parked at its final weld position.

\subsection{Discretization and time integration}

We discretize \cref{eq:heat} with continuous piecewise-linear ($P_1$) triangular
finite elements (\texttt{scikit-fem}~\cite{skfem2020}). The weak form yields the
semidiscrete system
\begin{equation}
  \bm{M}(T)\,\dot{\bm T} + \bm{A}(T)\,\bm T = \bm b(t),
\end{equation}
where $\bm M(T)$ is the (temperature-dependent) heat-capacity mass matrix with
entries $\int_\Omega \rho\,\cpapp(T)\,\phi_i\phi_j$, $\bm A=\bm K+\bm R$ collects the
conduction stiffness $\bm K=\int_\Omega k\,\grad\phi_i\cdot\grad\phi_j$ and the Robin
convection/radiation operator $\bm R$, and $\bm b$ is the Goldak load plus boundary
loads. Time integration uses the $\theta$-method (default backward Euler,
$\theta=1$),
\begin{equation}
  \Big(\tfrac{1}{\Delta t}\bm M + \theta\bm A\Big)\bm T^{n+1}
  = \tfrac{1}{\Delta t}\bm M\,\bm T^{n}
  - (1-\theta)\bm A\,\bm T^{n}
  + \theta\,\bm b^{n+1}+(1-\theta)\,\bm b^{n}.
\end{equation}
Because $\cpapp(T)$ (latent heat) and radiation make the problem nonlinear, each step
is solved by a Picard fixed-point iteration (re-assembling $\bm M$ and the linearized
radiation operator $\bm R_{\mathrm{rad}}=\varepsilon\sigma_{\mathrm{SB}}
(T^2+\Tinf^2)(T+\Tinf)$ at the current iterate) until the relative update falls below a
tolerance. Dirichlet constraints are imposed by condensation. The solver records a
time series of snapshots $\{T^n\}$ together with the torch frame and source state, and
serializes everything (including all process/boundary metadata) so the graph pipeline
can reconstruct every feature from a single self-contained file.

%==============================================================================
\section{Dataset generation}
\label{sec:dataset}
%==============================================================================

We generate a corpus designed so that a single generalization axis---the weld
\emph{trajectory topology}---is held out, while every other axis (geometry, plate
size, process parameters, boundary conditions) is covered by the training
distribution. This makes any generalization gap cleanly attributable to the
trajectory rather than to extrapolation in plate size or speed.

\paragraph{Sampling ranges.} Each simulation draws i.i.d.\ from:
plate bounding box \SIrange{0.08}{0.25}{m} per side; element edge $\approx$\SI{3}{mm}
(yielding $\sim$\numrange{800}{6600}-node meshes); gross power
\SIrange{1500}{3500}{W}; arc efficiency \numrange{0.70}{0.90}; travel speed
\SIrange{2}{12}{mm/s}; Goldak $b,c_f\in\SIrange{2.0}{3.5}{mm}$,
$c_r/c_f\in[1.6,2.4]$, $f_f\in[0.5,0.7]$; thickness \SIrange{4}{8}{mm}; ambient
\SIrange{290}{320}{K}; convection $h_{\mathrm{conv}}\in\SIrange{5}{40}{W.m^{-2}.K^{-1}}$;
emissivity in $\{0,0.2,0.35\}$ (weighted toward zero), with a 30\% chance that one outer
edge is held Dirichlet. Boundary-condition strata (radiation $\times$ Dirichlet) are
cycled to guarantee balanced coverage. Geometries are rectangular (50\%), L-shaped
(25\%), or holed (25\%); trajectories are straight, diagonal, sinusoidal, or
circular-arc paths.

\paragraph{Adaptive cooling tail.} After the weld phase, the source is switched off and
the simulation continues until the peak excess over ambient falls below 7\% of the
welding peak (near-field relaxed) or a cap of \SI{30}{s} is reached. This guarantees
every trajectory includes a complete local cool-down, which proved essential for
suppressing spurious far-field heating in long rollouts. The time step is pinned near
$\Delta t\approx\SI{0.05}{s}$ (matching the inference regime).

\paragraph{Corpus.} The final dataset is \num{360} simulations (\num{300} train /
\num{60} validation), \num{398597} snapshot-pair graphs in total. Splitting is done at
the \emph{simulation} level (every snapshot of a given weld lands in one fold) to
prevent temporal leakage of adjacent states across folds.

%==============================================================================
\section{Graph representation}
\label{sec:graph}
%==============================================================================

Each consecutive snapshot pair $(T^t,T^{t+1})$ of a simulation becomes one graph. The
mesh nodes are graph nodes; mesh edges (both orientations, no self-loops) are graph
edges. The learning target is the per-node temperature increment
$\bm y = T^{t+1}-T^{t}\in\Rset^{N}$.

A central design choice is to use \emph{only relative} features---no absolute
coordinates---so the model cannot memorize positions and must learn local physics:

\begin{description}[leftmargin=1.2em,style=nextline]
\item[Node features ($12$-d).]
$\big[\,T,\; q_{\mathrm{Goldak}},\; \xi',\;\eta',\; \eta_{\mathrm{eff}}P,\; v,\;
\mathbb{1}_{\mathrm{int}},\mathbb{1}_{\mathrm{Dir}},\mathbb{1}_{\mathrm{Rob}},\;
h_{\mathrm{conv}},\; \varepsilon,\; \Tinf\,\big]$.
Here $q_{\mathrm{Goldak}}$ is the analytical source \cref{eq:goldak} evaluated at the
node (gated to zero during the cooling tail); $(\xi',\eta')$ are the co-moving torch
coordinates in the $(\hat{\bm t},\hat{\bm n})$ frame; $\eta_{\mathrm{eff}}P$ and $v$ are
process scalars broadcast to all nodes; the three one-hots encode the node's
boundary type; and $h_{\mathrm{conv}},\varepsilon,\Tinf$ are the local boundary values.
The Goldak semi-axes are \emph{not} exposed as features: they are global per-simulation
constants already folded into $q_{\mathrm{Goldak}}$ and $(\xi',\eta')$.
\item[Edge features ($3$-d).] $[\,\Delta x,\Delta y,\norm{\bm u_{ij}}\,]$ with
$\bm u_{ij}=\bm x_i-\bm x_j$. Geometry enters only through these displacement vectors.
\end{description}

Continuous features and the target are standardized by z-score using statistics
computed over the training corpus; the three one-hot columns are left untouched.
Crucially, normalization constants are stored with the model so inference reproduces
the exact transform.

%==============================================================================
\section{MeshGraphNet architecture (the MPNN)}
\label{sec:mpnn}
%==============================================================================

The backbone is the Encoder--Processor--Decoder MeshGraphNet of
Pfaff et al.~\cite{pfaff2021learning}, a message-passing neural network (MPNN). Let
$\bm x_i\in\Rset^{12}$ be node features and $\bm e_{ij}\in\Rset^{3}$ edge features. All
sub-networks are multilayer perceptrons (MLPs) of width $d=128$ with two hidden layers,
ReLU activations, and a trailing LayerNorm (except the decoder).

\paragraph{Encoder.} Two independent MLPs lift raw features into a shared
$d$-dimensional latent space:
\begin{equation}
  \bm v_i^{(0)} = \mathrm{MLP}_{\mathrm{enc}}^{v}(\bm x_i),
  \qquad
  \bm \epsilon_{ij}^{(0)} = \mathrm{MLP}_{\mathrm{enc}}^{e}(\bm e_{ij}).
\end{equation}

\paragraph{Processor.} A stack of $L$ GraphNet blocks (here $L=8$). Each block updates
edges from their endpoints, aggregates messages at receiver nodes, and updates nodes,
all with \emph{residual} connections that let gradients survive the deep stack
(important for backpropagating through multi-step rollouts):
\begin{align}
  \bm\epsilon_{ij}^{(\ell+1)}
    &= \bm\epsilon_{ij}^{(\ell)}
      + \mathrm{MLP}_e^{(\ell)}\!\big([\,\bm\epsilon_{ij}^{(\ell)},\,
        \bm v_i^{(\ell)},\,\bm v_j^{(\ell)}\,]\big),
        \label{eq:edgeupdate}\\[2pt]
  \bm m_j^{(\ell+1)} &= \!\!\sum_{i\in\mathcal N(j)}\!\! \bm\epsilon_{ij}^{(\ell+1)},
        \qquad\text{(sum aggregation)}\\[2pt]
  \bm v_j^{(\ell+1)}
    &= \bm v_j^{(\ell)}
      + \mathrm{MLP}_v^{(\ell)}\!\big([\,\bm v_j^{(\ell)},\,\bm m_j^{(\ell+1)}\,]\big).
        \label{eq:nodeupdate}
\end{align}
With $L=8$ message-passing hops, information propagates up to 8 mesh edges per step,
giving the receptive field needed to represent the heat-affected zone around the torch.

\paragraph{Decoder.} A final MLP (no LayerNorm, so it can emit unbounded physical
values) projects node latents to the scalar output:
\begin{equation}
  \widetilde{\Delta T}_i = \mathrm{MLP}_{\mathrm{dec}}(\bm v_i^{(L)}) \in \Rset.
\end{equation}
In the \emph{baseline} (unconstrained) model this normalized increment is the
prediction. In the structure-preserving variants the decoder is bypassed and the
increment is built instead from the processor latents through the thermodynamic head
(\cref{sec:generic,sec:enthalpy}). The configuration used throughout has
$d=128$, $L=8$, $\approx$\num{1.29}M parameters for the unconstrained model.

\paragraph{Inference (autoregressive rollout).} At test time the model receives only
the initial field $T^0$ and the (known) Goldak schedule. Because the trajectory is
deterministic, all time-dependent features ($q_{\mathrm{Goldak}}$, $(\xi',\eta')$, $v$)
are precomputed for every step; during the loop only the temperature column is
overwritten with the running prediction:
\begin{equation}
  T^{t+1} = T^{t} + \mathcal{N}^{-1}_y\!\big(f_\theta(\bm x^t; \,\text{edges})\big),
\end{equation}
where $\mathcal N^{-1}_y$ de-normalizes the predicted $\Delta T$. The same feature
construction code is shared between training and rollout to guarantee parity.

%==============================================================================
\section{Structure-preserving thermodynamics: the GENERIC head}
\label{sec:generic}
%==============================================================================

The unconstrained MPNN is accurate but can violate physics over long rollouts
(spurious global heating, sub-ambient cooling, runaway). We therefore optionally
replace the free decoder increment by an increment that satisfies the laws of
thermodynamics \emph{by construction}, using the GENERIC
framework~\cite{grmela1997dynamics,ottinger2005beyond}.

\subsection{GENERIC in a nutshell}

GENERIC writes the evolution of a state $\bm z$ as a reversible plus a dissipative
part, driven by two potentials---total energy $E(\bm z)$ and entropy $S(\bm z)$:
\begin{equation}
  \frac{\dd \bm z}{\dd t}
  = \underbrace{\bm L(\bm z)\,\frac{\partial E}{\partial\bm z}}_{\text{reversible}}
  + \underbrace{\bm M(\bm z)\,\frac{\partial S}{\partial\bm z}}_{\text{dissipative}},
  \label{eq:generic}
\end{equation}
with $\bm L$ antisymmetric, $\bm M$ symmetric positive semi-definite (SPSD), and the
\emph{degeneracy conditions}
\begin{equation}
  \bm L\,\frac{\partial S}{\partial\bm z}=\bm 0,
  \qquad
  \bm M\,\frac{\partial E}{\partial\bm z}=\bm 0.
  \label{eq:degeneracy}
\end{equation}
These immediately yield the two laws:
\begin{equation}
  \frac{\dd E}{\dd t}=0
  \quad(\text{1st law}),
  \qquad
  \frac{\dd S}{\dd t}
  =\Big(\tfrac{\partial S}{\partial\bm z}\Big)^{\!\T}\!\bm M\,
   \Big(\tfrac{\partial S}{\partial\bm z}\Big)\ge 0
  \quad(\text{2nd law}).
  \label{eq:laws}
\end{equation}
Pure heat conduction is \emph{purely dissipative}: there is no reversible heat flow, so
$\bm L=\bm 0$ and $\dot{\bm z}=\bm M\,\partial S/\partial\bm z$ (plus a non-conservative
exchange with the external source/environment). Two choices then decide physical
consistency: \emph{what is the state $\bm z$}, and \emph{is
$\bm M\,\partial E/\partial\bm z=\bm0$ the correct energy conservation}.

\paragraph{Open, driven systems.}
The closed-form identities \cref{eq:laws} describe an \emph{isolated} system; welding is
\emph{open and driven}---the arc injects energy through the Goldak source and the
boundaries exchange heat by convection and radiation. We therefore use the standard
boundary/source extension of GENERIC for open systems~\cite{ottinger2005beyond}: the
structure-preserving operator acts only on the \emph{internal} bulk evolution, and the
exchange with the environment is added as a separate, non-conservative term
$\Delta^{\mathrm{ext}}$ evaluated after the internal dissipative
step~\cite{ottinger2006open,ottinger2005beyond}, in which boundary and bulk
contributions to the brackets are separated so that the isolated-system structure of
\cref{eq:laws} is recovered for the bulk while the boundary/source terms account for the
exchange with the environment,
\begin{equation}
  \dot{\bm z}=\underbrace{\bm M\,\partial S/\partial\bm z}_{\text{internal (degeneracy + 2nd law)}}
  \;+\;\underbrace{\bm f^{\mathrm{ext}}}_{\text{source + boundary fluxes}}.
  \label{eq:opengeneric}
\end{equation}
Consequently the degeneracy and second-law guarantees attach to the \emph{internal}
operator: $\bm 1^\T\bm M\,\partial S/\partial\bm z=0$ states that the irreversible bulk
dynamics conserves energy, while the total energy changes by \emph{exactly} the energy
crossing the boundary or deposited by the source (as physics demands, not a defect); and
$(\partial S/\partial\bm z)^\T\bm M\,(\partial S/\partial\bm z)\ge0$ is the non-negative
\emph{internal entropy production}, the correct second-law statement for an open system
(the system's total entropy may still fall as heat leaves through the boundary). The
external term $\bm f^{\mathrm{ext}}$ is realized by the (enriched) Goldak source and
Newton/Stefan cooling of \cref{sec:enriched}; it is the only part not constrained by the
GENERIC structure, by design.

\subsection{A learned SPSD operator on the graph}

We realize $\bm M$ as a \emph{graph Laplacian of learned conductances}. The processor
produces, for each undirected edge, a non-negative symmetric weight
\begin{equation}
  w_{ij} = e^{\alpha}\,\softplus\!\Big(\mathrm{MLP}\big(
    \bm v_i+\bm v_j,\;(\bm v_i-\bm v_j)^2,\;T_i+T_j,\;\abs{T_i-T_j}\big)\Big)
  = w_{ji}\ge 0,
  \label{eq:conductance}
\end{equation}
where $\bm v$ are processor node latents and $\alpha$ a learnable global log-scale
(initialized so the dissipation starts at $O(1\,\mathrm K)$ per step). Symmetry is
obtained \emph{for free} by feeding the MLP only features symmetric under
$i\leftrightarrow j$. Defining $\bm W=[w_{ij}]$, $\bm D=\operatorname{diag}(\bm W\bm 1)$,
the Laplacian is $\Lw=\bm D-\bm W$, which is SPSD and satisfies $\Lw\bm 1=\bm 0$ (zero
row/column sums).

The dissipative increment acts on the entropy gradient $\bmu$, the
\emph{thermodynamic force}:
\begin{equation}
  \mu_i \;\equiv\; \frac{\partial S}{\partial e_i} \;=\; \frac{1}{T_i},
  \label{eq:mu}
\end{equation}
which is the definition of (inverse) temperature, $\dd S=\dd E/T$, and holds for
\emph{any} $c_p(T)$, including latent heat. With $\bm M=\Lw$,
\begin{equation}
  \big(\Lw\bmu\big)_i = \sum_{j\sim i} w_{ij}\,(\mu_i-\mu_j).
  \label{eq:dissipation}
\end{equation}
A hot node ($\mu_i$ small) beside cold neighbors ($\mu_j$ large) receives a negative
increment---it \emph{loses} energy. Heat flows hot$\to$cold by construction.

\subsection{The inconsistency of a temperature state under latent heat}

It is tempting to take the state as temperature, $\bm z=\bm T$, with constant-$c_p$
potentials $E=\sum_i\rho\,c_p\,T_iV_i$ and $S=\sum_i\rho\,c_p\ln(T_i)V_i$, and to write
the temperature increment directly as $\Delta T_i=(\Lw\bmu)_i$. Then $\sum_i\Delta
T_i=\bm1^\T\Lw\bmu=0$ and one is tempted to call this ``energy conservation.'' It is
not. The thermodynamically correct dissipative law evolves \emph{energy}:
\begin{equation}
  \frac{\dd e_i}{\dd t}=(\Lw\bmu)_i,
  \qquad
  \frac{\dd T_i}{\dd t}=\frac{1}{\rho\,\cpapp(T_i)}\,\frac{\dd e_i}{\dd t}.
  \label{eq:correctT}
\end{equation}
Writing $\Delta T_i=(\Lw\bmu)_i$ directly \emph{drops} the factor
$1/(\rho\,\cpapp(T_i))$. When $c_p$ is constant this is a global rescale that the
learned $w_{ij}$ can absorb; but with latent heat $\cpapp$ varies by orders of
magnitude between cold metal and the melt pool, so
\begin{equation}
  \sum_i \Delta T_i = 0
  \;\;(\text{temperature-sum})
  \;\neq\;
  \sum_i \rho\,\cpapp(T_i)\,\Delta T_i = \sum_i \Delta e_i = 0
  \;\;(\text{energy}).
  \label{eq:wrongcons}
\end{equation}
The two differ precisely by the $\cpapp(T)$ factor that blows up at the melt pool.
Physically, \cref{eq:correctT} says $\dd T/\dd t$ should be \emph{small} there (incoming
energy is swallowed by the phase change); the temperature-state head has no such damping
and over-moves the hottest nodes. This explains the empirical observation that the
temperature-state structure-preserving head \emph{caps/distorts the peak} while the
unconstrained network fits it better.

%==============================================================================
\section{The enthalpy-state GENERIC head}
\label{sec:enthalpy}
%==============================================================================

The fix is the classical \emph{enthalpy method} of phase-change heat
transfer~\cite{voller1987enthalpy}: take energy (enthalpy) as the conserved state and
recover temperature by inverting the latent-heat-aware curve.

\subsection{State, potentials, and the two laws}

Let the proper conserved state be the \emph{extensive nodal energy}
$\bm z=\bm{\mathcal E}$, $\mathcal E_i = h_i\,V_i$ [\si{J}], where $h_i=h(T_i)$ is the
volumetric enthalpy (\cref{eq:enthalpy}, monotone hence invertible) and $V_i$ is the
node's tributary volume (the diagonal of the lumped FEM mass matrix). The total energy
is then $E=\sum_i\mathcal E_i$, so $\partial E/\partial\mathcal E_i = 1$ and the GENERIC
degeneracy
\begin{equation}
  \bm M\,\frac{\partial E}{\partial\bm z}=\Lw\bm 1=\bm 0
  \label{eq:degen_exact}
\end{equation}
holds \emph{exactly on any mesh} (the graph Laplacian has zero column sums). The entropy
gradient is $\mu_i=\partial S/\partial\mathcal E_i=1/T_i$ as before (the definition of
temperature, valid for any $c_p$). The dissipative evolution and temperature recovery are
\begin{equation}
  \boxed{\;\frac{\dd \mathcal E_i}{\dd t} = (\Lw\bmu)_i
    = \sum_{j\sim i} w_{ij}\Big(\frac{1}{T_i}-\frac{1}{T_j}\Big),
    \qquad T_i = h^{-1}\!\big(\mathcal E_i/V_i\big).\;}
  \label{eq:enthalpyevol}
\end{equation}

\paragraph{Remark (implemented uniform-volume reduction).}
Choosing the \emph{volumetric} enthalpy $\bm h$ as the state (rather than the extensive
$\bm{\mathcal E}$) gives $\partial E/\partial h_i=V_i$, so the degeneracy would require
$\Lw\bm V=\bm 0$ rather than $\Lw\bm 1=\bm 0$; with a standard graph Laplacian these
coincide \emph{only} for uniform $V_i$. The meshes used here are uniform tensor-product
triangulations (\cref{sec:dataset}), so we take equal interior tributary volumes
$V_i\equiv V$: the constant $V$ is absorbed into the learned conductance scale
$e^{\alpha}$, the head evolves the intensive temperature-equivalent enthalpy $H$
(\cref{sec:conditioning}) directly, and conservation reduces to the unweighted
$\sum_i\Delta H_i^{\mathrm{diss}}=0$. This is identical to the extensive form
\cref{eq:enthalpyevol} in the mesh interior and differs only at boundary/corner nodes
(tributary volume $V/2$, $V/4$), a small residual confined to $\partial\Omega$. The
general $V_i$-weighted form \cref{eq:enthalpyevol} restores \emph{exact} conservation on
arbitrary (e.g.\ arc-refined) meshes at the cost of a single learning-free weighting;
we adopt it in the theory and flag the weighting as an implementation refinement
(\cref{sec:discussion}).

\begin{proposition}[Discrete first and second laws, any mesh]
\label{prop1}
For any non-negative symmetric conductances $w_{ij}=w_{ji}\ge 0$ and any nodal volumes
$V_i>0$, the dissipative update \cref{eq:enthalpyevol} satisfies, per graph,
\begin{equation}
  \frac{\dd E}{\dd t}=\sum_i \frac{\dd \mathcal E_i}{\dd t}
    = \bm 1^\T \Lw\bmu = 0
  \quad\text{(energy conserved)},
  \qquad
  \frac{\dd S}{\dd t}=\bmu^\T\Lw\bmu
  = \tfrac12\sum_{i,j} w_{ij}(\mu_i-\mu_j)^2 \ge 0.
\end{equation}
\end{proposition}
\begin{proof}
$\bm1^\T\Lw=\bm0$ because $\Lw$ has zero column sums, giving the first law; and $\Lw$ is
SPSD because $w_{ij}\ge0$ and symmetric, so the entropy-production quadratic form is
non-negative. Crucially, the conserved quantity is now the \emph{extensive energy}
$\mathcal E=\sum_i\mathcal E_i$ (not the temperature sum, nor the unweighted volumetric
enthalpy sum), so this is the genuine first law on \emph{any} mesh, even with latent
heat. The uniform-volume implementation (above) conserves $\sum_i\Delta
H_i^{\mathrm{diss}}=0$, which equals $\dd E/\dd t=0$ exactly when $V_i$ is constant (mesh
interior) and up to the $O(V/2,\,V/4)$ boundary-node correction otherwise.
\end{proof}

Latent heat is now handled exactly: at the melt pool, $T=h^{-1}(h)$ is nearly flat
(the plateau in \cref{fig:hcurve}), so a large $\Delta h$ produces almost no $\Delta T$,
just as in the FEM.

\subsection{Connection to Fourier conduction}

To check that \cref{eq:enthalpyevol} reduces to Fourier conduction, note
$\grad T=-T^2\grad(1/T)=-T^2\grad\mu$, so the conductive flux is
$\bm q=-k\grad T=k\,T^2\grad\mu$. Discretizing $\grad\!\cdot\!\bm q$ on the graph yields
exactly the Laplacian form $\Lw\bmu$ with edge conductances scaling as
$w_{ij}\sim k\,T_iT_j/\ell_{ij}^2$. The network's task is to learn these
(temperature- and geometry-dependent) conductances.

\subsection{Numerical conditioning: temperature-equivalent enthalpy}
\label{sec:conditioning}

Raw volumetric enthalpy is $h\sim\SI{e9}{J.m^{-3}}$, so the inverse map has slope
$\partial T/\partial h=1/(\rho\cpapp)\sim\num{2e-7}$. This factor multiplies every
gradient flowing back to the conductances and gains, pushing them below Adam's
$\varepsilon=\num{e-8}$ floor and stalling training (empirically: the loss freezes and
diffusion never develops). We therefore work in \emph{temperature-equivalent enthalpy}
(units of kelvin):
\begin{equation}
  H(T) \;=\; \frac{h(T)}{\rho\,\cpsens}
  \;=\; \int_{T_0}^{T}\frac{\cpapp(\tau)}{\cpsens}\,\dd\tau .
  \label{eq:Heq}
\end{equation}
Away from phase change $H\approx T-T_0$ (unit slope); across melting/boiling the slope
spikes (the plateau). Since this is a \emph{constant} rescale of $h$, energy conservation
is unchanged ($\sum_i\Delta H_i^{\mathrm{diss}}=0\Leftrightarrow\sum_i\Delta
h_i^{\mathrm{diss}}=0$), but now increments are $O(\mathrm K)$ and gradients $O(1)$.
This single rescaling moved the conductance gradients from $\sim$\num{e-13} to
$\sim$\num{e-4} and unstuck training from the first epoch.

\subsection{Discrete scheme}

The curve $H(T)$ is tabulated once from the \emph{same} material model used by the FEM,
on a uniform grid $T^{(0)}<\dots<T^{(G-1)}$ ($G=4096$ from \SIrange{200}{8000}{K}):
\begin{equation}
  H^{(g)} = \sum_{m\le g}\frac{\cpapp(T^{(m)})}{\cpsens}\,\Delta T^{(m)} ,
\end{equation}
which is strictly increasing (hence invertible) since $\cpapp>0$. The arrays
$(T^{(g)},H^{(g)})$ ride in the model's \texttt{state\_dict} so checkpoints reproduce the
exact physics. Both $H=\mathcal H(T)$ and $T=\mathcal H^{-1}(H)$ are evaluated by a
differentiable piecewise-linear interpolation (binary search for the bracket, linear
blend), whose derivative is the local slope---finite and well defined, and inside the
autograd graph. One forward step per node:
\begin{align}
  \Delta H_i^{\mathrm{diss}}
    &= \sum_{j\sim i} w_{ij}\,(\mu_i-\mu_j),
       \qquad \mu_i=1/T_i,
       \label{eq:dHdiss}\\
  H_i^{\mathrm{new}}
    &= \mathcal H(T_i) + \Delta H_i^{\mathrm{diss}} + \Delta H_i^{\mathrm{ext}},
       \label{eq:Hupdate}\\
  \Delta T_i
    &= \mathcal H^{-1}\!\big(H_i^{\mathrm{new}}\big) - T_i,
       \label{eq:dT}
\end{align}
and the head returns the normalized $\Delta T_i$, so the rollout loop and data pipeline
are unchanged (the network still rolls temperature).

%==============================================================================
\section{The enriched source}
\label{sec:enriched}
%==============================================================================

Both structure-preserving heads (temperature and enthalpy) were observed to run
\emph{cold at the active peak}. The reason is structural: the dissipative operator can
\emph{only cool} the hottest node (hot$\to$cold by construction); the \emph{only} term
that can heat the peak is the external source. With a single global scalar gain it was
too rigid to inject enough energy. The external exchange is, however, exactly the
\emph{non-conservative} term that the GENERIC degeneracy and second law do \emph{not}
constrain---so we may give it expressivity without breaking any guarantee.

We replace the scalar source/cooling gains by per-node learned modulations:
\begin{equation}
  \Delta H_i^{\mathrm{ext}}
  = \underbrace{\softplus\!\big(g_q+\mathrm{MLP}_q(\bm v_i,q_i,T_i)\big)\,q_i}
    _{\text{Goldak heating }\ge 0}
  + \underbrace{\softplus\!\big(g_c+\mathrm{MLP}_c(\bm v_i,T_i,\Tinf)\big)\,
    r_i\,(\Tinf-T_i)}_{\text{Newton cooling on Robin nodes }r_i},
  \label{eq:enriched}
\end{equation}
where $q_i$ is the node's Goldak source feature and $r_i$ the Robin one-hot. The MLP
output layers are zero-initialized, so the head \emph{starts identical} to the scalar
version and then learns where to inject more energy. The heating gain is
$\softplus(\cdot)\ge0$ and is gated by $q_i$ (it cannot invent far-field energy); the
cooling term keeps the correct sign. The dissipation \cref{eq:dHdiss} remains fully
structure-preserving. The gains carry a $\Delta t/(\rho c_p)$-type scale, initialized so
$\softplus(g_q)\approx\SI{0.025}{}$ (a per-step deposit of $q\Delta t$ yields the right
$\sim$\SI{e2}{K} peak heating) and $\softplus(g_c)\approx\num{e-3}$.

This enriched-source enthalpy GENERIC model has \num{1375238} parameters
($+\,$\num{83461} over the unconstrained backbone: the conductance MLP, the two source
MLPs, and the global scales).

%==============================================================================
\section{Training}
\label{sec:training}
%==============================================================================

\subsection{Loss, optimizer, and training noise}

The base objective is the mean squared error on the normalized increment
$\Delta T=T^{t+1}-T^{t}$. We optimize with Adam (learning rate \num{e-3}, weight decay
\num{e-4}, gradient-norm clipping at $1.0$) and a plateau learning-rate scheduler. To
stabilize autoregressive rollouts we inject zero-mean Gaussian noise into the input
temperature during training and correct the target accordingly (the standard
MeshGraphNet noise of Pfaff et al.~\cite{pfaff2021learning}):
\begin{equation}
  \tilde T_i^{t}=T_i^{t}+\eta_i,\quad \eta_i\sim\mathcal N(0,\sigma^2);
  \qquad
  \Delta T_i^{\mathrm{target}}=T_i^{t+1}-\tilde T_i^{t}=y_i-\eta_i,
  \label{eq:noise}
\end{equation}
so the network learns to drive a perturbed state back onto the true trajectory. We use
$\sigma=\SI{5}{K}$. Noise is applied only to the training fold and only to the (raw)
temperature column, before normalization.

\subsection{Push-forward (multi-step) training}
\label{sec:pushforward}

A recurring empirical lesson was that improving the single-step objective harder (lower
training MSE, lower short-rollout validation RMSE) did \emph{not} improve, and sometimes
\emph{worsened}, the long inference rollout: the single-step objective is simply not
aligned with the \num{1526}-step inference horizon. We therefore train with a
\emph{push-forward} objective: unroll the model $K$ steps and backpropagate through the
whole unroll (full backpropagation through time).

Concretely, the data loader yields windows of $K$ consecutive same-simulation graphs
$[g_t,\dots,g_{t+K-1}]$. We seed temperature from the (noised) true field at the window
start and roll forward, feeding the model's own prediction back as the next input's
temperature column, accumulating a scaled-temperature MSE over the window
(\cref{alg:pf}). $K=1$ recovers the classic single-step objective. We use $K=2$, which
makes the operator contractive over two-step composition---enough, combined with the
thermodynamic structure, to align with the long rollout while remaining cheap.

\begin{algorithm}[t]
\caption{Push-forward training step (window of $K$ step-aligned batches)}
\label{alg:pf}
\begin{algorithmic}[1]
\State \textbf{input:} batches $[B_0,\dots,B_{K-1}]$, normalizer $\mathcal N$, model $f_\theta$
\State $\hat T \gets$ \Call{DynamicNoise}{$B_0.T,\,B_0.\Tinf$} \Comment{seed from window start}
\State $\mathcal L \gets 0$
\For{$k=0$ \textbf{to} $K-1$}
  \State $\bm x \gets B_k.\bm x$ with temperature column $\gets \hat T$
  \State $\Delta T \gets \mathcal N_y^{-1}\big(f_\theta(\text{standardize}(\bm x), B_k.\text{edges})\big)$
  \State $T_{\mathrm{pred}} \gets \hat T + \Delta T$
  \State $T_{\mathrm{true}} \gets B_k.T + B_k.y$
  \State $\mathcal L \gets \mathcal L + \big\lVert (T_{\mathrm{pred}}-T_{\mathrm{true}})/\sigma_T\big\rVert^2 / N$
  \State $\hat T \gets T_{\mathrm{pred}}$ \Comment{feed forward; \emph{no} detach $\Rightarrow$ BPTT}
\EndFor
\State $\mathcal L \gets \mathcal L / K$; \quad backprop, clip, optimizer step
\end{algorithmic}
\end{algorithm}

\subsection{Validation metric and model selection}

Validation uses the \emph{full autoregressive rollout} RMSE on held-out simulations
(not single-step error), averaged over the held-out fold; the best checkpoint is the one
minimizing this metric. We additionally evaluate every checkpoint on the out-of-
distribution spiral (below), because even the rollout validation does not perfectly
track the very-long inference horizon.

\subsection{Configuration and compute}

The headline run uses $d=128$, $L=8$ message-passing steps, batch size 16, $K=2$
push-forward, $\sigma=\SI{5}{K}$ noise, plateau scheduler, 80 epochs, simulation-level
split (\num{288} train / \num{72} val, validation rollout capped at 20 simulations for
cost). Training was performed on a single NVIDIA A100-80GB ($\sim$\SI{15}{min}/epoch,
$\sim$\SI{58}{GB} VRAM at this batch/$K$); the processed-graph cache lives on a network
volume so re-launches skip the build. Local evaluation/rollouts run on an RTX-class GPU.

%==============================================================================
\section{Experiments and results}
\label{sec:results}
%==============================================================================

\subsection{Out-of-distribution test: the Archimedean spiral}

The held-out generalization probe is a constant-speed Archimedean-spiral weld
$\bm\gamma(\theta)=\bm c+b\,\theta\,[\cos\theta,\sin\theta]$ on a square plate of side
\SI{0.2}{m}, 3 turns, travel speed \SI{10}{mm/s}, net power $\eta P$ with
$P=\SI{2500}{W}$, $\eta=0.8$. This trajectory is \emph{never} seen in training (all
training paths are straight/diagonal/sinusoid/arc), yet plate size and speed are
interior to the training ranges (\cref{sec:dataset}), so the spiral isolates trajectory
generalization. The FEM ground truth has \num{4225} nodes and \num{1526} time steps
(\SI{76.25}{s})---an autoregressive horizon many times longer than any training window.
The spiral additionally exercises genuine \emph{re-heating} physics: successive spiral
arms re-heat one another, a multi-pass effect absent from the single-pass training data.

\subsection{Headline result}

The enriched-source enthalpy GENERIC model trained with push-forward $K=2$ attains:
\begin{itemize}[leftmargin=1.4em]
  \item \textbf{Held-out single-pass validation rollout RMSE} that is \emph{stable}
  across epochs and does not degrade (best \SI{34.5}{K} at epoch 60), unlike the
  $K=1$ variant which improved validation while drifting the spiral worse;
  \item \textbf{Spiral global RMSE \SI{33.2}{K}} (best-validation checkpoint), about
  $2\times$ better than the best non-structure-preserving model obtained at any point
  in the study (\SI{65}{K}), and \emph{stable} over all \num{1526} steps (bounded
  $\max\abs{\text{err}}$, no divergence);
  \item \textbf{Maximum-principle respect}: minimum temperature stays within
  $\sim$\SI{6}{K} of the physical floor (ambient/Robin bath), never cooling
  sub-ambient.
\end{itemize}
\Cref{fig:spiral} shows the five-probe thermal histories (FEM vs.\ surrogate) at radii
spanning the melt pool to the cool corner. The shapes and arrival times of the passing
peaks are captured; the one remaining flaw is a mild smoothing of the sharpest central
peak (\SI{1713}{K} vs.\ FEM \SI{1921}{K}), while the mid-field is essentially nailed
(e.g.\ \SI{908}{K} vs.\ \SI{949}{K}).

\figIfExists{spiral_pf_best_ep60_thermal_history.png}{0.95}{%
  Spiral out-of-distribution test: FEM ground truth vs.\ the enriched-source enthalpy
  GENERIC surrogate (push-forward $K{=}2$, best-validation checkpoint), at five probe
  locations from the hot center to the cool corner. Global RMSE \SI{33.2}{K} over
  \num{1526} autoregressive steps; stable, no divergence.}{fig:spiral}

\subsection{Ablation: thermodynamic structure ON vs.\ OFF}

To isolate the value of the structure, we train an otherwise identical
\emph{unconstrained} MeshGraphNet (GENERIC OFF, decoder active) with the same
push-forward $K=2$, the same split, budget, noise, and hyperparameters. The two models
therefore differ \emph{only} in the head, and their validation metrics are directly
comparable. We then sweep every checkpoint (epochs 10--80) on the spiral and measure not
just best-case RMSE but the \emph{stability signals}: rollout-RMSE variance across
checkpoints, the minimum temperature (maximum-principle violation), and whether the
deployable best-validation checkpoint is also the best on the spiral.

\Cref{tab:ablation} summarizes the ablation; \cref{fig:ablation} plots spiral-RMSE and
minimum-temperature versus epoch for both models. The findings:
\begin{itemize}[leftmargin=1.4em]
  \item \textbf{Maximum principle.} The structured model respects the physical floor
  (min $T\approx\SI{287}{K}$, $\sim$\SI{6}{K} undershoot); the unconstrained model cools
  to \SIrange{140}{190}{K} at \emph{every} checkpoint---up to $\sim$\SI{150}{K} of
  unphysical sub-ambient cooling.
  \item \textbf{Variance.} The structured model's spiral RMSE sits in a tight band
  (33--61\,K, $\sigma\approx3.4$); the unconstrained model swings 38--111\,K
  ($\sigma\approx10.2$, $\sim$$3\times$ larger).
  \item \textbf{Selectability.} For the structured model the best-validation checkpoint
  \emph{is} the best spiral checkpoint (validation tracks the long rollout); for the
  unconstrained model the best-validation checkpoint gives a mediocre spiral
  (\SI{58.6}{K}) and the best spiral checkpoint cannot be selected by validation.
  \item \textbf{Held-out single-pass.} The structured model is also better on ordinary
  held-out rollouts (e.g.\ \SI{8.1}{K} vs.\ \SI{36.2}{K} on one validation simulation).
\end{itemize}

\begin{table}[t]
\centering
\caption{Ablation: GENERIC enthalpy structure ON vs.\ OFF, both push-forward $K{=}2$,
identical otherwise. The structured model respects the maximum principle, has
$\sim$$3\times$ lower spiral-RMSE variance, and a validation metric that tracks the
spiral.}
\label{tab:ablation}
\small
\begin{tabular}{lccccc}
\toprule
Configuration & Params & Val.\ RMSE & Spiral RMSE & Spiral range & Min.\ $T$ \\
              &        & (K)        & @best-val (K) & ep10--80 [$\sigma$] (K) & / undershoot (K)\\
\midrule
\textbf{ON} (enthalpy GENERIC      & \multirow{2}{*}{\num{1375238}} & \multirow{2}{*}{34.5}
  & \multirow{2}{*}{33.2} & \multirow{2}{*}{33--61 [3.4]} & \multirow{2}{*}{287 ($-6$)}\\
\;\;+ enriched source) & & & & &\\
\addlinespace[2pt]
\textbf{OFF} (pure MeshGraphNet)   & \num{1291777} & 46.9 & 58.6 & 38--111 [10.2] & 140 ($-153$)\\
\bottomrule
\end{tabular}
\par\vspace{2pt}
{\footnotesize Physical floor $=\SI{293}{K}$ (ambient $=$ Robin bath $=$ initial).}
\end{table}

\figIfExists{ablation_figure.png}{0.95}{%
  Ablation (both push-forward $K{=}2$). \emph{Left:} spiral global RMSE vs.\ epoch; the
  structured model (ON) stays in a tight band and its best-validation checkpoint
  ($\star$) coincides with its best spiral, whereas the unconstrained model (OFF) swings
  $\sim$$3\times$ more and its best-validation checkpoint is mediocre on the spiral.
  \emph{Right:} minimum temperature vs.\ epoch; ON stays at the physical floor
  (\SI{293}{K} line), OFF cools sub-ambient at every checkpoint.}{fig:ablation}

\figIfExists{compare_spiral_fem_on_off.png}{0.95}{%
  Three-way spiral thermal histories (five probes): FEM (black), GENERIC ON (green),
  pure MeshGraphNet OFF (red, dashed). The unconstrained model's far-field and
  cool-corner behavior is visibly less faithful; per-curve RMSE in the legend.}%
  {fig:comparespiral}

\subsection{Held-out single-pass validation}

\Cref{fig:val1,fig:val2} show held-out single-pass rollouts (in-distribution
trajectories, unseen geometry/process/boundary). The structured model is near-perfect
in the far field (RMSE $\sim$\SI{1}{K}, no spurious echo) and accurate at the peak; the
residual error concentrates only in extreme dynamic-range outliers. These results---not
the OOD spiral---constitute the core deployable claim: strong generalization across
geometry, plate size, process parameters, and boundary conditions.

\figIfExists{compare_val_sim001_fem_on_off.png}{0.92}{%
  Held-out single-pass validation simulation (four peak-fraction probes): FEM (black),
  GENERIC ON (green), OFF (red, dashed). The structured model tracks all probes; global
  RMSE \SI{8.1}{K} (ON) vs.\ \SI{36.2}{K} (OFF).}{fig:val1}

\figIfExists{compare_val_sim006_fem_on_off.png}{0.92}{%
  A second held-out validation simulation. Global RMSE \SI{36.0}{K} (ON) vs.\
  \SI{63.4}{K} (OFF); the ON residual is concentrated at the sharp peak, the far field
  is near-exact.}{fig:val2}

%==============================================================================
\section{Discussion}
\label{sec:discussion}
%==============================================================================

\paragraph{What the structure buys.} The honest framing of the ablation is not that the
unconstrained network is catastrophic---at a lucky checkpoint it reaches a respectable
number---but that it is \emph{fragile}: you cannot select its good checkpoint by any
deployable metric, it violates the maximum principle everywhere, and it varies
$\sim$$3\times$ more across training. The GENERIC structure instead provides
\emph{horizon-independent guarantees}: energy conservation and $\dd S/\dd t\ge0$ hold at
every single step regardless of rollout length, which is exactly what a \num{1526}-step
(or longer) deployment needs. This is a stronger and more defensible claim than a small
RMSE improvement.

\paragraph{Where the work is done.} A caveat we make explicit: because the enriched
source is a learned per-node MLP (gated by the Goldak field), one may ask whether the
flexible source is ``painting'' the temperature rather than the structured dissipation
doing the physics. The source is constrained (non-negative heating, gated by $q$,
sign-correct cooling), but a principled next step is to replace the source MLP by a
\emph{named} temperature-dependent arc-efficiency $\eta(T)$ with few parameters: if a
bounded physical mechanism still recovers the peak, the effect is real physics rather
than black-box flexibility.

\paragraph{Limitations and future improvements.}
(i) \textbf{Exact energy conservation vs.\ the implemented uniform-volume reduction.}
The first law in \cref{prop1} holds on \emph{any} mesh for the extensive-energy state
$\mathcal E_i=h_iV_i$ \cref{eq:enthalpyevol}. The deployed head uses the uniform-volume
reduction $V_i\equiv V$, which is exact in the interior of the uniform tensor meshes used
here and leaves only a small residual at boundary/corner nodes (tributary volumes
$V/2,\,V/4$). On strongly non-uniform meshes---e.g.\ refined near the arc, as is common in
production welding FEM---the unweighted conservation $\sum_i\Delta h_i=0$ would no longer
equal the physical $\sum_i V_i\Delta h_i=0$. We emphasize that this does \emph{not} affect
the reported results: on the present uniform meshes the two formulations are numerically
indistinguishable away from $\partial\Omega$, so the uniform-volume head is an excellent
approximation, and the maximum-principle/stability guarantees observed empirically are
unaffected. Restoring \emph{exact} conservation on arbitrary meshes requires only the
cheap, learning-free volume weighting (the lumped FEM mass diagonal, already available);
we leave the $V_i$-weighted formulation and its retraining on refined meshes as a
straightforward future improvement.
(ii) The model is 2D (thickness-integrated); a 3D extension follows the same graph
machinery.
(iii) The constant-$\rho c_p$ rescaling in \cref{sec:conditioning} is exact only up to the
curve $H(T)$ used; the curve is the FEM's, so the surrogate inherits the FEM's material
model.
(iv) The training corpus has a small soft-cap tail of very hot ($>$\SI{4000}{K})
simulations; tightening the vaporization cap and regenerating would clean it (the enthalpy
head's $H(T)$ table auto-syncs).
(v) The OOD spiral's central peak is still mildly smoothed; sharper-peak modeling is
future work.

%==============================================================================
\section{Related work}
%==============================================================================

\emph{Learned mesh simulators.} MeshGraphNets~\cite{pfaff2021learning} and graph network
simulators~\cite{sanchez2020learning} established message passing on simulation meshes;
our backbone is a configurable MeshGraphNet with the standard training-noise trick for
rollout stability. \emph{Structure-preserving learning.} The GENERIC
framework~\cite{grmela1997dynamics,ottinger2005beyond} has been used to learn
thermodynamically consistent dynamics (e.g.\ GENERIC neural networks and structure-
preserving operators)~\cite{hernandez2021structure}; our contribution is to identify
that imposing it on \emph{temperature} is inconsistent under latent heat and to fix it
with the classical \emph{enthalpy method}~\cite{voller1987enthalpy}, realized as a
learned SPSD graph-Laplacian on a graph-network surrogate. \emph{Welding/AM thermal
surrogates} have used POD/ROM and CNN/GNN models; the novelty here is the combination of
thermodynamic structure, phase-change-consistent energy bookkeeping, and push-forward
training for long-horizon stability.

%==============================================================================
\section{Conclusion}
%==============================================================================

We presented a MeshGraphNet surrogate for transient welding heat transfer that is
\emph{thermodynamically consistent by construction}. The key technical insight is that
GENERIC structure must be imposed on \emph{energy}, not temperature: with latent heat,
conserving the sum of nodal temperatures is not conserving energy, and the mismatch is
largest exactly at the melt pool. Carrying the volumetric enthalpy as the conserved
state (the enthalpy method), evolving it with a learned SPSD graph-Laplacian, and
recovering temperature through the latent-heat-aware curve restores genuine energy
conservation and the second law through phase change. An enriched, non-conservative
source recovers melt-pool peak heating without breaking the guarantees, and push-forward
$K{=}2$ training aligns the objective with long autoregressive rollouts. On a held-out
out-of-distribution spiral weld of \num{1526} steps the model reaches \SI{33.2}{K} global
RMSE while respecting the maximum principle and remaining stable, and a matched ablation
shows the structure delivers low-variance, selectable, physically admissible rollouts
where the unconstrained network is accurate-but-fragile.

%==============================================================================
\section*{Reproducibility}
%==============================================================================
\addcontentsline{toc}{section}{Reproducibility}
The pipeline comprises: the FEM solver (\texttt{thermal\_solver.py}: Goldak source,
apparent heat capacity, $\theta$-method + Picard), dataset generation
(\texttt{generate\_dataset.py}: stratified sampling and adaptive cooling), the graph
builder (\texttt{graph\_builder.py}: 12-d relative node features, 3-d edge features), the
model (\texttt{meshgraphnet.py}: Encoder--Processor--Decoder plus the
\texttt{EnthalpyGenericThermalHead}), and training (\texttt{train.py}: push-forward
unroll, training noise, simulation-level split). The structure-preserving guarantees
($\sum_i\Delta H^{\mathrm{diss}}_i=0$, $\dd S/\dd t\ge0$, latent-heat damping) are
unit-tested. Normalization constants and the enthalpy table ride in the model
\texttt{state\_dict}, so checkpoints reproduce the exact physics.

%==============================================================================
\begin{thebibliography}{9}
%==============================================================================
\bibitem{pfaff2021learning}
T.~Pfaff, M.~Fortunato, A.~Sanchez-Gonzalez, and P.~W.~Battaglia.
\newblock Learning mesh-based simulation with graph networks.
\newblock In \emph{ICLR}, 2021.

\bibitem{sanchez2020learning}
A.~Sanchez-Gonzalez, J.~Godwin, T.~Pfaff, R.~Ying, J.~Leskovec, and P.~W.~Battaglia.
\newblock Learning to simulate complex physics with graph networks.
\newblock In \emph{ICML}, 2020.

\bibitem{grmela1997dynamics}
M.~Grmela and H.~C.~\"Ottinger.
\newblock Dynamics and thermodynamics of complex fluids. I. Development of a general formalism.
\newblock \emph{Physical Review E}, 56(6):6620--6632, 1997.

\bibitem{ottinger2005beyond}
H.~C.~\"Ottinger.
\newblock \emph{Beyond Equilibrium Thermodynamics}.
\newblock Wiley, 2005.

\bibitem{ottinger2006open}
H.~C.~\"Ottinger.
\newblock Nonequilibrium thermodynamics for open systems.
\newblock \emph{Physical Review E}, 73(3):036126, 2006.

\bibitem{goldak1984new}
J.~Goldak, A.~Chakravarti, and M.~Bibby.
\newblock A new finite element model for welding heat sources.
\newblock \emph{Metallurgical Transactions B}, 15(2):299--305, 1984.

\bibitem{voller1987enthalpy}
V.~R.~Voller and C.~Prakash.
\newblock A fixed grid numerical modelling methodology for convection-diffusion mushy
region phase-change problems.
\newblock \emph{Int.\ J.\ Heat Mass Transfer}, 30(8):1709--1719, 1987.

\bibitem{hernandez2021structure}
Q.~Hern\'andez, A.~Bad\'ias, D.~Gonz\'alez, F.~Chinesta, and E.~Cueto.
\newblock Structure-preserving neural networks.
\newblock \emph{Journal of Computational Physics}, 426:109950, 2021.

\bibitem{skfem2020}
T.~Gustafsson and G.~D.~McBain.
\newblock scikit-fem: A Python package for finite element assembly.
\newblock \emph{Journal of Open Source Software}, 5(52):2369, 2020.
\end{thebibliography}

\end{document}
